(1) 速度平方为 \(\ell^2[\dot\theta^2+\sin^2\theta\,\dot\phi^2]\)，势能取 \(U=mg\ell(1-\cos\theta)\)。拉格朗日量 \(L=\frac12m\ell^2(\dot\theta^2+\sin^2\theta\,\dot\phi^2)-U\)。循环变量 \(\phi\) 的守恒动量为 \(L_z=m\ell^2\sin^2\theta\,\dot\phi\)，从而 \(\dot\phi=L_z/(m\ell^2\sin^2\theta)\)。把方位动能用固定角动量重写，得一维有效势 \(U_{\rm eff}=mg\ell(1-\cos\theta)+L_z^2/[2m\ell^2\sin^2\theta]\)。

(2) 有效势的驻点条件 \(mg\ell\sin\theta-L_z^2\cos\theta/[m\ell^2\sin^3\theta]=0\)，即 \(L_z^2=m^2g\ell^3\sin^4\theta/\cos\theta\)。由于左侧正，必须 \(\cos\theta>0\)，故匀速圆锥态只有 \(0<\theta<\pi/2\)。定义角速度 \(\Omega=\dot\phi\)，上述条件化为 \(\Omega^2=g/(\ell\cos\theta)\)，相当于水平方向向心加速度与竖向重力的同时平衡；对固定非零 \(L_z\)，右侧函数 \(\sin^4\theta/\cos\theta\) 从零单调增加至无穷，所以平衡角唯一。

(3) 在固定 \(L_z\) 下计算有效势二阶导数 \(U_{\rm eff}''=mg\ell\cos\theta+(L_z^2/(m\ell^2))[1/\sin^2\theta+3\cos^2\theta/\sin^4\theta]\)。在平衡点代入 \(L_z^2/(m^2\ell^3)=g\sin^4\theta_e/\cos\theta_e\)，可化为 \(U_{\rm eff}''=mg\ell[\cos\theta_e+(\sin^2\theta_e+3\cos^2\theta_e)/\cos\theta_e]\)，最终 \(U_{\rm eff}''=mg\ell(1+3\cos^2\theta_e)/\cos\theta_e\)。极角惯量为 \(m\ell^2\)，故 \(\omega_\theta^2=g(1+3\cos^2\theta_e)/(\ell\cos\theta_e)>0\)，证明圆锥运动对小极角扰动稳定。

(4) 若 \(\theta_e\to0\)，圆周半径趋零，角速度 \(\Omega^2\to g/\ell\)，而极角模式 \(\omega_\theta\to2\sqrt{g/\ell}\)；当 \(\theta_e\to\pi/2^-\)，\(\Omega\) 与 \(\omega_\theta\) 同时发散，说明要维持近水平圆锥需要极大转速与张力。竖向平衡给细线张力 \(T\cos\theta_e=mg\)，所以 \(T=mg/\cos\theta_e\)。张力必须为正，若球摆允许到达绳松弛或支点碰撞状态，本约束模型需另行修正。
